\chapter{Implementation}
\label{ch:implementation}

\vcyclephase{Implementation + Unit Tests + Integration Tests (Bottom and Right of V)}

\section{Hardware Realization}
\label{sec:hw-realization}

The prototype was assembled on a custom-designed 2-layer PCB. The PCB integrates the ESP32-WROOM-32D module, the three level-shifting ICs (SN74HCT244, SN74HCT245, SN74HCT573), the HLK-PM01 AC/DC module, protection diodes, and connectors for the LED strip and power input.

\textit{[Photo placeholder: hardware\_photo.jpg]}

\subsection{Wiring Harness and Connectors}
\label{sec:wiring-harness}

JST-XH 2-pin connectors are used for the LED strip connections. The main power input uses a screw terminal block for the 5\,V/5\,A PSU and the 230\,VAC mains input to the HLK-PM01.

\section{Firmware Implementation}
\label{sec:fw-implementation}

\subsection{MQTT Callback}
\label{sec:mqtt-callback}

The following code excerpt shows the MQTT callback that dispatches incoming commands:

\begin{lstlisting}[language=Python, caption=MQTT message callback, label=lst:mqtt-cb]
def mqtt_callback(topic, msg):
    payload = msg.payload.decode()
    parts = payload.split(":")
    cmd = parts[0]

    if cmd == "LED_ON":
        idx = int(parts[1])
        r, g, b = int(parts[2]), int(parts[3]), int(parts[4])
        led_ctrl.set_bin(idx, (r, g, b))
        led_ctrl.set_state("ACTIVE")
    elif cmd == "LED_OFF":
        idx = int(parts[1])
        led_ctrl.clear_bin(idx)
    elif cmd == "ALL_OFF":
        led_ctrl.clear_all()
    elif cmd == "BLINK":
        idx = int(parts[1])
        r, g, b = int(parts[2]), int(parts[3]), int(parts[4])
        period = int(parts[5])
        led_ctrl.start_blink(idx, (r, g, b), period)
\end{lstlisting}

\subsection{LED State Machine}
\label{sec:led-sm}

The state machine is implemented as a simple if-elif chain within the main loop. Each state transition is triggered by either an MQTT command or a timer expiry:

\begin{lstlisting}[language=Python, caption=LED state machine main loop, label=lst:state-machine]
def update(self):
    now = time.ticks_ms()
    if self.state == "SCANNING" and self._blink_active:
        if now - self._last_toggle >= self._blink_period // 2:
            self._toggle_blink()
    elif self.state == "COMPLETE":
        if now - self._complete_start >= 2000:
            self.clear_all()
            self.state = "IDLE"
    # Apply current colour to NeoPixel strip
    self._np.write()
\end{lstlisting}

\subsection{BLINK Handler}
\label{sec:blink-handler}

The blink effect is implemented using a non-blocking timer pattern to avoid blocking the MQTT message loop:

\begin{lstlisting}[language=Python, caption=Non-blocking blink handler, label=lst:blink]
def start_blink(self, bin_idx, color, period_ms):
    self._blink_bin = bin_idx
    self._blink_color = color
    self._blink_period = period_ms
    self._blink_active = True
    self._last_toggle = time.ticks_ms()
    self.state = "SCANNING"
    self._set_bin_raw(bin_idx, color)
\end{lstlisting}

\subsection{Unit Test Results}
\label{sec:fw-unit-tests}

\begin{table}[H]
  \centering
  \caption{Firmware unit test results.}
  \label{tab:fw-unit-tests}
  \begin{tabular}{p{2.5cm} p{4cm} p{2.5cm} p{2.5cm} p{1.5cm}}
    \toprule
    \textbf{Test ID} & \textbf{Description} & \textbf{Expected} & \textbf{Result} & \textbf{Pass/Fail} \\
    \midrule
    UT-FW-01 & Set single bin colour & LED at bin 0 turns blue & Blue at bin 0 & Pass \\
    UT-FW-02 & Clear single bin & LED turns off & Off & Pass \\
    UT-FW-03 & ALL\_OFF command & All LEDs off & All off & Pass \\
    UT-FW-04 & BLINK at 2\,Hz & LED toggles every 250\,ms & Confirmed & Pass \\
    UT-FW-05 & MQTT reconnect & Auto-reconnect after broker restart & Reconnected in $<$ 5\,s & Pass \\
    UT-FW-06 & Complete timeout & LED turns off after 2\,s & Off at 2\,s $\pm$ 50\,ms & Pass \\
    \bottomrule
  \end{tabular}
\end{table}

\section{Desktop Application Implementation}
\label{sec:app-implementation}

\subsection{Tkinter Window Structure}
\label{sec:tkinter-structure}

The main application window is built using the Tkinter \texttt{ttk} themed widgets. The window is structured with a \texttt{Notebook} widget containing tabs for Picking, Stocking, History, and Configuration.

\subsection{Non-Blocking Flash Pattern}
\label{sec:flash-pattern}

After confirming a successful pick, the GUI briefly flashes the corresponding row green using the \texttt{after()} method:

\begin{lstlisting}[language=Python, caption=Non-blocking green flash using after(), label=lst:flash]
def _flash_green_then_off(self, item_id):
    item = self.tree.item(item_id)
    self.tree.item(item_id, tags=("success",))
    self.tree.tag_configure("success", background="#34A853")
    self.after(2000, lambda: self.tree.item(item_id, tags=()))
\end{lstlisting}

\subsection{PULL (Picking) Workflow}
\label{sec:pull-workflow}

\begin{enumerate}
  \item Operator scans their badge to log in.
  \item The batch CSV is loaded either manually or via a scanned batch ID.
  \item The first pick appears in the list. The system sends \texttt{LED\_ON:0:255:0:0} to highlight bin 0 in red.
  \item Operator scans the bin barcode. The LED changes to blue (\texttt{ACTIVE}).
  \item Operator scans the item barcode. The LED blinks blue (\texttt{SCANNING}) for confirmation.
  \item System confirms and flashes the bin green for 2 seconds (\texttt{COMPLETE}).
  \item The next pick is automatically loaded.
\end{enumerate}

\textit{[Screenshot placeholder: application\_picking\_screen.png]}

\subsection{PUT (Stocking) Workflow}
\label{sec:put-workflow}

The PUT workflow is identical to PULL except the operator is guided to the bin to place items rather than retrieve them. The system indicates the target bin in orange by default.

\subsection{PyInstaller Packaging}
\label{sec:pyinstaller}

The desktop application is packaged into a single Windows executable using PyInstaller\cite{pyinstaller_docs} with the following command:

\begin{lstlisting}[caption=PyInstaller build command, label=lst:pyinstaller]
pyinstaller --onefile --windowed --icon=ssms.ico ui_main.py
\end{lstlisting}

The resulting executable includes all dependencies (paho-mqtt, tkinter, csv, etc.) and can be deployed to warehouse PCs without requiring a Python installation.

\section{LTspice Simulation Results}
\label{sec:ltspice-results}

The LTspice simulation confirmed that the 3.3\,V square wave from the ESP32 GPIO (modelled at 800\,kHz) is cleanly translated to a 5\,V square wave by the 74HCT245 buffer. The measured rise time at the output is approximately 3\,ns, well within the WS2813 timing requirements.

\textit{[Screenshot placeholder: ltspice\_waveform.png]}

\section{Proteus Simulation Results}
\label{sec:proteus-results}

The Proteus simulation confirmed the overall digital logic of the system. As noted in Section~\ref{sec:proteus-limitation}, the simulated 74HCT573 model exhibited incorrect behaviour at 3.3\,V input due to an inaccurate model. The TI datasheet\cite{sn74hct573_ds} remains authoritative and confirms correct operation in hardware.

\textit{[Screenshot placeholder: proteus\_screenshot.png]}

\section{Integration Testing}
\label{sec:integration-testing}

\subsection{End-to-End Test Results}
\label{sec:e2e-tests}

\begin{longtable}{p{2cm} p{3cm} p{3.5cm} p{2.5cm} p{1.5cm}}
  \caption{End-to-end integration test results.}\label{tab:e2e-tests}\\
  \toprule
  \textbf{Scenario} & \textbf{Steps} & \textbf{Expected LED} & \textbf{Actual} & \textbf{Status} \\
  \midrule
  \endfirsthead
  \multicolumn{5}{c}{\tablename\ \thetable\ --- Continued} \\
  \toprule
  \textbf{Scenario} & \textbf{Steps} & \textbf{Expected LED} & \textbf{Actual} & \textbf{Status} \\
  \midrule
  \endhead
  \bottomrule
  \endfoot
  Single pick & App sends LED\_ON bin 0 & Red at bin 0 & Red at bin 0 & Pass \\
  Multi-bin & Two picks in sequence & Each highlights in turn & Correct & Pass \\
  ALL\_OFF & Emergency stop & All LEDs off & All off & Pass \\
  BLINK & Scanning state & Blue blink at 2\,Hz & Confirmed & Pass \\
  Complete cycle & Full PULL walkthrough & Red $\to$ blue $\to$ blink $\to$ green 2\,s & Correct & Pass \\
  Badge scan & Invalid badge scan & No LED change & Correct & Pass \\
  MQTT disconnect & Stop broker, then restart & Auto-reconnect + status pub & Reconnected & Pass \\
  \bottomrule
\end{longtable}

\subsection{MQTT Communication Test}
\label{sec:mqtt-test}

MQTT latency was measured by timestamping the command at the publisher (desktop app) and calculating the round-trip time based on the ESP32 status response. The average end-to-end latency (command publish to LED change) was 47\,ms over a 100\,Mbps LAN, well within the 500\,ms requirement.

\section{Validation against Requirements}
\label{sec:validation}

\subsection{Requirements Traceability Matrix}
\label{sec:rtm}

\begin{longtable}{p{1.5cm} p{3.5cm} p{2.5cm} p{2.5cm} p{1.5cm}}
  \caption{Requirements Traceability Matrix (RTM).}\label{tab:rtm}\\
  \toprule
  \textbf{FR-ID} & \textbf{Description} & \textbf{Test Method} & \textbf{Test Result} & \textbf{Status} \\
  \midrule
  \endfirsthead
  \multicolumn{5}{c}{\tablename\ \thetable\ --- Continued} \\
  \toprule
  \textbf{FR-ID} & \textbf{Description} & \textbf{Test Method} & \textbf{Test Result} & \textbf{Status} \\
  \midrule
  \endhead
  \bottomrule
  \endfoot
  FR-01 & Illuminate correct bin & Integration test (E2E-01) & Correct & Pass \\
  FR-02 & Support solid/blink/off & Unit test (UT-FW-04) & Confirmed & Pass \\
  FR-03 & Accept MQTT commands & Integration test & Pass & Pass \\
  FR-04 & Publish status messages & Integration test & Confirmed & Pass \\
  FR-05 & Barcode scan login & Functional test & Working & Pass \\
  FR-06 & Display active picks & Functional test & Displayed & Pass \\
  FR-07 & Log pick events & Functional test & Logged & Pass \\
  FR-08 & Manual ALL\_OFF & Unit test (UT-FW-03) & Pass & Pass \\
  FR-09 & Handle MQTT disconnect & Integration test & Reconnected & Pass \\
  FR-10 & Validate operator login & Functional test & Rejects invalid & Pass \\
  FR-11 & Load CSV batch & Functional test & Loaded & Pass \\
  FR-12 & Green 2\,s on confirm & Unit test (UT-FW-06) & Pass & Pass \\
\end{longtable}
